\subsection{Common Projective Origin of Quantum and Gravitational Structure}
  \label{subsec:common-projective-origin}

  Sections~\ref{subsec:relation-to-general-relativity}
  and~\ref{subsec:interpretation-of-quantum-phenomena} treat the recovery of
  general relativity and quantum mechanics separately.
  It is worth making explicit that both emerge from the same structural source:
  the non-injectivity of~$\Pi$ and the Born--Infeld admissibility bound,
  acting on the same static substrate~$\chi$.

  On the quantum side, non-injectivity implies that any observable~$o$ carries
  an irreducible fiber $\Pi^{-1}(o)$ of indistinguishable pre-images.
  The effective description, having no access to~$\chi$ directly, must represent
  the coexistence of projectively admissible contributions internally.
  This forces a complex, norm-preserving, additive structure---the Hilbert space
  framework---and, in the regime where the Born--Infeld bound keeps fluctuations
  subcritical and the background remains quasi-stable, a unique effective dynamics
  of Schr\"odinger type (Section~\ref{sec:relation-to-quantum-formalism},
  Appendix~\ref{app:schrodinger-derivation}).
  Unitarity is not postulated: it holds if and only if finite projective resolution
  and bounded admissible flux hold simultaneously
  (Theorem~\ref{thm:schrodinger-emergence}, condition~C5).

  On the gravitational side, the same non-injectivity implies that the projection
  $\Pi$ cannot be globally injective, so that the effective metric $g_{\mu\nu}$
  is not a primitive object but a descriptive construct encoding variations in
  the local projective ordering rate
  (Section~\ref{subsec:emergent-curvature}).
  Matter configurations locally saturate projective capacity, producing
  differential ordering rates that are reinterpreted as effective curvature.
  The Born--Infeld bound, which on the quantum side enforces unitarity, on the
  gravitational side selects the unique infrared dynamics compatible with
  descriptive invariance and prevents the emergence of arbitrary higher-derivative
  corrections (Section~\ref{sec:necessity-noninjective}).

  The two recoveries are therefore not independent achievements of the framework.
  They are two aspects of a single structural fact: a static substrate~$\chi$
  accessible only through a non-injective projection~$\Pi$ bounded by a
  Born--Infeld saturation constraint cannot produce an effective description that
  is simultaneously faithful (injective), linear (quantum), and geometric (relativistic).
  Each of the three adjectives describes the same compression from a different
  observational vantage point.
  The apparent conceptual distance between quantum mechanics and general relativity
  is a consequence of the fact that they are two effective languages for the same
  underlying projective structure, optimised for complementary regimes of
  resolution and scale.
